\documentclass[12pt,a4paper]{article}
\usepackage[utf8]{utf8}
\usepackage[margin=1in]{geometry}
\usepackage{hyperref}
\usepackage{xcolor}
\usepackage{listings}
\usepackage{booktabs}
\usepackage{amsmath}
\usepackage{array}
\usepackage{enumitem}

% Colors
\definecolor{primary}{RGB}{124, 58, 237}
\definecolor{darkbg}{RGB}{15, 23, 42}
\definecolor{codebg}{RGB}{241, 245, 249}
\definecolor{codegreen}{RGB}{16, 185, 129}
\definecolor{codered}{RGB}{239, 68, 68}

\hypersetup{
    colorlinks=true,
    linkcolor=primary,
    filecolor=primary,      
    urlcolor=primary,
    pdftitle={HOI Multi-Provider Lounge Access System Architecture Master Plan},
    pdfpagemode=FullScreen,
}

\lstset{
    backgroundcolor=\color{codebg},
    basicstyle=\ttfamily\small,
    breaklines=true,
    captionpos=b,
    commentstyle=\color{gray},
    keywordstyle=\color{primary}\bfseries,
    stringstyle=\color{codegreen},
    showstringspaces=false,
    tabsize=2,
    frame=single
}

\title{\textbf{\Large HOI Multi-Provider Lounge Access System}\\[0.5em] \Large System Architecture \& Integration Master Plan}
\author{\textbf{Engineering \& Architecture Team}}
\date{\today}

\begin{document}

\maketitle
\tableofcontents
\newpage

% ============================================================
\section{Executive Summary}
% ============================================================
This document provides the complete, authoritative specification and implementation blueprint for the \textbf{HOI Multi-Provider Lounge Access Control System}.

The system bridges IoT hardware scanners (\textbf{ESP32 Turnstile Controllers}) with multiple independent 3rd-party provider backends (e.g., HOI LoungePass, DreamFolks, Encalm) via a centralized, high-performance \textbf{Middleware Proxy Relay Server}.

Key highlights of this architectural specification include:
\begin{itemize}[noitemsep]
    \item \textbf{Hardware-agnostic 3rd Party Routing}: Seamless integration of 8-character alphanumeric QR tokens without modifying scanner hardware code.
    \item \textbf{Dual-Layer Security Model}: Hardware-level HMAC-SHA256 device authentication combined with Provider-level Bearer JWT Validation.
    \item \textbf{Granular Access Control}: Built-in support for QR Expiry Limits, Usage Count Restrictions (Single-use vs Multi-use), and real-time audit logging.
    \item \textbf{Zero-Downtime Provider Migration}: Database-driven device mapping that allows remote provider switching without re-flashing ESP32 firmware.
\end{itemize}

% ============================================================
\section{High-Level System Architecture}
% ============================================================

The architecture consists of three decoupled operational tiers:

\begin{enumerate}
    \item \textbf{Tier 1: ESP32 Hardware Edge Scanners (Gate Controllers)}: Scans 8-character QR tokens at gate turnstiles and transmits cryptographically signed payloads to the Middleware Server.
    \item \textbf{Tier 2: Main Middleware Relay Server (The Proxy Router)}: Verifies hardware integrity via HMAC, queries device-to-provider mappings from the relational database, and proxies HTTP requests to the target provider.
    \item \textbf{Tier 3: 3rd Party Provider Servers (HOI / Partners)}: Validates QR tokens against business rules (Expiry, Max Usage, Registration Status) and returns approval or rejection codes.
\end{enumerate}

\begin{lstlisting}[caption={System Architecture Data Flow Diagram}]
[ ESP32 Gate Scanner ]
         |
         |-- 1. Scans 8-char QR ("Z5NUFKDY")
         |-- 2. Signs payload with HMAC-SHA256 (deviceId|timestamp|token)
         |-- 3. POST request to Middleware Server (/api/esp32/scan)
         v
[ YOUR MAIN MIDDLEWARE SERVER ]  (The Proxy Router)
         |
         |-- 4. Authenticate ESP32 HMAC Signature (Hardware Auth)
         |-- 5. Query DB `devices` table for `device_id`
         |      Returns: Provider="HOI", Target_URL="https://hoi-api.com"
         |-- 6. Forward / Relay QR Token to Provider API Endpoint:
         |      POST https://hoi-api.com/lounge-visits/enquiries/Z5NUFKDY
         |      Header: Authorization: Bearer <HOI_BEARER_TOKEN>
         v
[ HOI 3rd PARTY SERVER ]
         |
         |-- 7. Validates Bearer Token, Expiry & Max Usage Count
         |-- 8. Updates status to "COMPLETED" & Returns HTTP Response
         v
[ YOUR MAIN MIDDLEWARE SERVER ]
         |
         |-- 9. Log transaction in DB `scan_logs`
         |-- 10. Transmit final action to ESP32 (GATE_OPEN / GATE_DENIED)
         v
[ ESP32 Gate Scanner ]
         |-- 11. Triggers Relay (Turnstile Opens! Green LED Bleeps!)
\end{lstlisting}

\newpage
% ============================================================
\section{Component Specifications}
% ============================================================

\subsection{Component A: ESP32 Hardware Edge Scanner}
The ESP32 microcontroller is connected to an 8-character QR optical scanner module and a turnstile relay switch.

\subsubsection{Request Payload Format}
When a QR code is scanned, the ESP32 calculates an HMAC-SHA256 signature using its factory-provisioned \texttt{secretKey} and posts to the Middleware:

\begin{lstlisting}[language=json, caption={ESP32 to Middleware Request Payload}]
{
    "deviceId": "ESP_DEL_T3_01",
    "timestamp": 1788624000,
    "scannedToken": "Z5NUFKDY",
    "signature": "e3b0c44298fc1c149afbf4c8996fb92427ae41e4649b934ca495991b7852b855"
}
\end{lstlisting}

\subsubsection{Signature Formula}
\[
\text{Payload} = \text{deviceId} \parallel \text{"|"} \parallel \text{timestamp} \parallel \text{"|"} \parallel \text{scannedToken}
\]
\[
\text{Signature} = \text{HMAC-SHA256}(\text{Payload}, \text{secretKey})
\]

\subsection{Component B: Main Middleware Server (Proxy Router)}
The Middleware Server is responsible for zero-latency routing and auditing:
\begin{enumerate}
    \item \textbf{Hardware Integrity}: Re-computes the HMAC-SHA256 signature to guarantee the request originated from an authorized ESP32 hardware device.
    \item \textbf{Device Lookup}: Queries the \texttt{devices} database table using the \texttt{deviceId} to retrieve the target Provider's API URL and pre-configured Bearer Token.
    \item \textbf{HTTP Proxying}: Dispatches an asynchronous HTTP POST request to the Provider's endpoint.
    \item \textbf{Gate Command Translation}: Converts Provider HTTP status codes into binary hardware commands (\texttt{GATE\_OPEN} vs \texttt{GATE\_DENIED}).
\end{enumerate}

\subsection{Component C: 3rd Party Provider Server (HOI Mock API)}
The 3rd Party Provider validates the business logic:
\begin{itemize}
    \item \textbf{JWT Bearer Authorization}: Validates caller token.
    \item \textbf{Record Existence}: Ensures the 8-character Enquiry ID exists in storage (Returns \texttt{HTTP 404 Not Found} if invalid).
    \item \textbf{Token Expiry}: Checks \texttt{expires\_at} timestamp (Returns \texttt{HTTP 410 Gone} if expired).
    \item \textbf{Usage Count Control}: Increments \texttt{used\_count} against \texttt{max\_usage} limit (Returns \texttt{HTTP 409 Conflict} if limit exceeded).
\end{itemize}

\newpage
% ============================================================
\section{Database Schema Design (DDL)}
% ============================================================

The system utilizes two primary relational database tables: \texttt{devices} (Configuration) and \texttt{scan\_logs} (Auditing).

\subsection{Table 1: \texttt{devices} (Hardware-to-Provider Mapping)}

\begin{lstlisting}[language=SQL, caption={DDL Statement for devices Table}]
CREATE TABLE devices (
    id INT AUTO_INCREMENT PRIMARY KEY,
    device_id VARCHAR(50) NOT NULL UNIQUE,
    device_name VARCHAR(100) NOT NULL,
    location_name VARCHAR(150) NOT NULL,
    provider_name VARCHAR(50) NOT NULL,
    provider_api_url VARCHAR(255) NOT NULL,
    provider_bearer_token TEXT NOT NULL,
    secret_key VARCHAR(128) NOT NULL,
    status ENUM('ACTIVE', 'INACTIVE') DEFAULT 'ACTIVE',
    created_at TIMESTAMP DEFAULT CURRENT_TIMESTAMP,
    updated_at TIMESTAMP DEFAULT CURRENT_TIMESTAMP ON UPDATE CURRENT_TIMESTAMP
) ENGINE=InnoDB DEFAULT CHARSET=utf8mb4;
\end{lstlisting}

\subsection{Table 2: \texttt{scan\_logs} (Audit History \& Analytics)}

\begin{lstlisting}[language=SQL, caption={DDL Statement for scan\_logs Table}]
CREATE TABLE scan_logs (
    id INT AUTO_INCREMENT PRIMARY KEY,
    device_id VARCHAR(50) NOT NULL,
    scanned_token VARCHAR(20) NOT NULL,
    provider_name VARCHAR(50) NOT NULL,
    http_code INT NOT NULL,
    gate_action ENUM('GATE_OPEN', 'GATE_DENIED') NOT NULL,
    response_msg TEXT,
    scanned_at TIMESTAMP DEFAULT CURRENT_TIMESTAMP,
    INDEX idx_device (device_id),
    INDEX idx_scanned_at (scanned_at)
) ENGINE=InnoDB DEFAULT CHARSET=utf8mb4;
\end{lstlisting}

% ============================================================
\section{Error Handling \& Gate Action Matrix}
% ============================================================

The following matrix specifies the exact behavior of the Middleware Server and ESP32 Edge Device for every possible HTTP response status from the 3rd-party provider:

\begin{table}[h!]
\centering
\small
\begin{tabular}{|m{2.2cm}|m{2.5cm}|m{4.5cm}|m{3.5cm}|}
\hline
\textbf{Provider Status} & \textbf{Error Code} & \textbf{Middleware Action} & \textbf{ESP32 Hardware Action} \\ \hline
\textbf{HTTP 200 OK} & \texttt{COMPLETED} & Logs success; Returns \texttt{GATE\_OPEN} & 🚪 Turnstile Opens + Green LED + Chime \\ \hline
\textbf{HTTP 401} & \texttt{unauthorized} & Logs auth failure; Returns \texttt{GATE\_DENIED} & 🔴 Red LED + Display "Auth Error" \\ \hline
\textbf{HTTP 404} & \texttt{enquiry\_not\_found} & Logs invalid ID attempt; Returns \texttt{GATE\_DENIED} & 🔴 Red LED + Display "Invalid QR ID" \\ \hline
\textbf{HTTP 409} & \texttt{limit\_exceeded} & Logs limit exceeded; Returns \texttt{GATE\_DENIED} & 🔴 Red LED + Display "Already Used" \\ \hline
\textbf{HTTP 410} & \texttt{qr\_token\_expired} & Logs QR expiry; Returns \texttt{GATE\_DENIED} & 🔴 Red LED + Display "QR Expired" \\ \hline
\textbf{HTTP 500 / 504} & \texttt{server\_error} & Logs timeout; Returns \texttt{GATE\_DENIED} & 🔴 Red LED + Display "Network Timeout" \\ \hline
\end{tabular}
\caption{Complete Gate Action Matrix}
\end{table}

\newpage
% ============================================================
\section{API Specifications}
% ============================================================

\subsection{1. Obtain OAuth 2.0 Access Token}
\begin{itemize}
    \item \textbf{Method}: \texttt{POST}
    \item \textbf{Endpoint}: \texttt{/oauth/token}
    \item \textbf{Content-Type}: \texttt{application/json}
\end{itemize}

\begin{lstlisting}[language=json, caption={OAuth 2.0 Token Request & Response}]
// Request Body
{
    "grant_type": "client_credentials",
    "client_id": "client_hoi_prod",
    "client_secret": "secret_hoi_lounge_2024_key"
}

// Response Body (HTTP 200 OK)
{
    "access_token": "eyJhbGciOiJIUzI1NiIsInR5cCI6IkpXVCJ9...",
    "token_type": "Bearer",
    "expires_in": 2592000,
    "scope": "lounge_read_write"
}
\end{lstlisting}

\subsection{2. Validate Lounge Visit / Enquiry ID}
\begin{itemize}
    \item \textbf{Method}: \texttt{POST}
    \item \textbf{Endpoint}: \texttt{/lounge-visits/enquiries/\{ENQUIRY\_ID\}}
    \item \textbf{Header}: \texttt{Authorization: Bearer <ACCESS\_TOKEN>}
\end{itemize}

\begin{lstlisting}[language=json, caption={Lounge Visit Validation Response}]
// Success Response (HTTP 200 OK)
{
    "statusCode": 200,
    "success": true,
    "message": "Lounge visit processed & completed successfully!",
    "data": {
        "enquiry_id": "Z5NUFKDY",
        "passenger_name": "Rahul Sharma",
        "status": "COMPLETED",
        "max_usage": 1,
        "used_count": 1,
        "remaining_uses": 0,
        "expires_at": "2026-10-05 23:59:59"
    }
}
\end{lstlisting}

% ============================================================
\section{Security Audit \& Compliance Standards}
% ============================================================
\begin{enumerate}
    \item \textbf{RFC 6749 Compliance}: Implements standard OAuth 2.0 Client Credentials Grant for B2B API integrations.
    \item \textbf{FIPS 198-1 HMAC Compliance}: Uses SHA-256 keyed-hash algorithm to prevent hardware spoofing and replay attacks.
    \item \textbf{Timezone Standardization}: All logs, timestamps (\texttt{created\_at}, \texttt{expires\_at}), and audit trails are synchronized to \textbf{Indian Standard Time (IST - Asia/Kolkata, UTC+5:30)}.
\end{enumerate}

\end{document}
